\documentstyle[%


righttag%


,11pt%


,amssymb%


,russian%               remove in Eng.


,izvru%                 Set izven% in Eng.


]{amsart}





%


\newcommand{\re}{\operatorname{Re}}


\newcommand{\im}{\operatorname{Im}}


\newcommand{\pr}{\operatorname{pr}}


\newcommand{\tts}[1]{}





%\hoffset=-15mm%                  For Eng.


\setcounter{page}{21}


\god{2001}


\nomer{5~(468)} %                        Remove in Eng.


%\nomer{Vol.\,45, No.\,5}%               Set in Eng.


%\str{\pageref{begin}--\pageref{end}}%  Set in Eng.


\udk{517.54}





\author{\it В.Н.\,Дубинин, В.Ю.\,Ким }


% NB:   ^^ remove in Eng.


\title{Усредняющее преобразование множеств


и функций\\ на римановых поверхностях}


\thanks{Pабота выполнена при поддержке Российского фонда


фундаментальных исследований (грант 99-01-00443) и


программы ``Университеты России'' (грант 991282).}





\date{}


%\pagestyle{myheadings}%         Set in Eng.


\pagestyle{plain} %              Remove in Eng.


\begin{document}


%\thispagestyle{plain}%          Set in Eng.


\maketitle


\label{begin}





Усредняющее преобразование областей, введенное Маркусом в [1], получило


значительное применение в геометрической теории функций комплексного


переменного (напр., [2]--[4]). Отметим работы [2], [3], в


которых, в частности, даны обобщения этого преобразования и многочисленные


приложения к различным классам аналитических функций. В данной статье


преобразование Маркуса [1] развивается в ином направлении. Усредняются


множества, лежащие на римановых поверхностях над комплексной плоскостью с


учетом их листности накрытия. Аналогичные преобразования для круговой


симметризации осуществлялись в [5]--[8]


и других работах. В идейном плане в данной работе


происходит соединение усреднения [1] и кусочно-разделяющей симметризации [4].


Однако построенные здесь преобразования не сводятся к ним. В качестве


приложений даются новые теоремы о покрытии отрезков, о совместно $p$-листных


областях и неравенствах для рациональных функций.





\section{Линейно-усредняющее преобразование}





Пусть $\cal R$ --- риманова поверхность, лежащая над комплексной плоскостью


$z=x+iy$, $\cal R (\Pi)$ --- множество точек поверхности $\cal R$, проекция


которых принадлежит полосе $\Pi=\{z: 0\le x\le 1\}$. Говоря об открытых и


замкнутых множествах на $\cal R(\Pi)$, будем иметь в виду относительную


топологию, являющуюся следом на $\cal R(\Pi)$ топологии, заданной в $\cal R$.


Совокупность открытых множеств $\{\cal D\}$ назовем допустимым семейством


множеств, если выполняются следующие условия: каждое множество ${\cal D}$


принадлежит множеству $\cal R(\Pi)$ некоторой римановой поверхности $\cal R$;


множества $\{\cal D\}$ совместно $p$-листные, т.\,е. каждая точка $z\in\Pi$


покрывается не более чем $p$ различными точками, принадлежащими множествам из


совокупности $\{{\cal D}\}$ (с учетом кратности); существует число $c>0$ такое,


что множества $\{\cal D\}$\; $p$-кратно покрывают множество $\Pi_c=\{z: 0\le


x\leq 1,\ |y|\geq c\}$, т.\,е. над каждой точкой $z\in\Pi_c$ лежит ровно $p$


точек с учетом кратности из множеств $\{{\cal D}\}$; наконец, ни одно из


множеств совокупности $\{{\cal D}\}$ не содержит жордановой дуги, целиком


покрывающей прямую $\re z=x$ при каком-либо $x$, $0<x<1$. Далее, семейство


замкнутых множеств $\{\cal E\}$ назовем допустимым семейством, соответствующим


семейству $\{{\cal D}\}$, если каждое множество ${\cal E}$ из $\{{\cal E}\}$


принадлежит некоторому множеству $\cal D \in \{\cal D\}$ и если множества


совокупности $\{\cal E\}$


также $p$-кратно покрывают полуполосы $\Pi_c$ при некотором $c>0$.


Для простоты изложения удобно считать, что каждое множество $\cal D$ семейства


$\{{\cal D}\}$ содержит некоторое замкнутое множество ${\cal E} \equiv \cal


E({\cal D})$ семейства $\{{\cal E}\}$, возможно, пустое. Обозначим через


$l_c(x,A)$ линейную меру Лебега множества точек из $A$, лежащих над отрезком


$[x-ic,\ x+ic]$. Результатом линейно-усредняющего преобразования ${\cal L}^p$


допустимого семейства открытых множеств $\{{\cal D}\}$ назовем множество


${\cal L}^p \{{\cal D}\}=\{z=x+iy\in\Pi: y>c-\frac{1}{2p}\sum\limits_{\{\cal


D\}}l_c(x,{\cal D})\}$, где $c$ из определения семейства


$\{{\cal D}\}$. Легко видеть, что ${\cal L}^p\{{\cal D}\}$ не зависит от


выбора числа $c>0$. Аналогично определяется линейно-усредняющее преобразование


допустимого семейства замкнутых множеств


${\cal L}^p\{{\cal E}\}=


\{z=x+iy\in\Pi: y\geq c-\frac{1}{2p}\sum\limits_{\{\cal


E\}}l_c(x,{\cal E})\}$.





Предположим, что вещественная функция $\omega(Z)$ задана на множествах


${\cal D}$ допустимого семейства открытых множеств $\{{\cal D}\}$ и


удовлетворяет следующим условиям:


\begin{itemize}


\item[1)] $0\le \omega (Z)\leq 1$


на множествах ${\cal D}$ и $\omega (Z)=1$ только на


множествах ${\cal E}$ некоторого допустимого семейства замкнутых множеств


$\{{\cal E}\}$, соответствующего семейству $\{{\cal D}\}$;


\item[2)] $\omega (Z)$ непрерывная на ${\cal D}$ и гармоническая в $\cal


D\setminus{\cal E}$, ${\cal D}\in \{{\cal D}\}$, ${\cal E}\in\{{\cal E}\}$;


\item[3)] при любом фиксированном $\alpha$, $ 0<\alpha<1$, множества


${\cal D}(\alpha)=\{Z\in{\cal D}: \omega(Z)>\alpha\}$,


${\cal D}\in\{{\cal D}\}$,


образуют совокупность из конечного числа непересекающихся областей,


имеющих разве лишь конечное число точек ветвления;


\item[4)] $\lim\limits_{Z\to\partial{\cal D}}\omega(Z)=0$,


${\cal D}\in\{{\cal D}\}$.


\end{itemize}


Определим линейно-усредняющее преобразование функции $\omega(Z)$


как переход от этой функции к функции ${\cal L}^p\omega(z)$, заданной в


полуполосе $\Pi^+=\{z\in\Pi: y\geq 0\}$ формулой


$$


{\cal L}^p\omega(x+iy)=


\begin{cases}


1, &y\ge c-\frac{1}{2p}\sum\limits_{\{{\cal E}\}}l_c(x,{\cal E});\\


\alpha, &y=c-\frac{1}{2p}\sum\limits_{\{{\cal D}\}}l_c(x,{\cal D}(\alpha));\\


0, &y\leq c-\frac{1}{2p}\sum\limits_{\{{\cal D}\}}l_c(x,{\cal D}),\\


\end{cases}


$$


где $c$ --- некоторое число из определения семейств $\{{\cal D}\}$ и


$\{{\cal E}\}$. Для произвольной функции $v$, заданной на некоторой римановой


поверхности ${{\cal R}}$, обозначим через


$I(v,\Omega)=\int\limits_{\Omega}|\nabla v|^2$


интеграл Дирихле от функции $v$ по множеству $\Omega \subset {\cal R}$


(напр., [9], с.\,256).





\begin{thm}


Пусть $\{{\cal D}\}$ --- допустимое семейство открытых множеств, а


$\{{\cal E}\}$ --- соответствующее ему семейство замкнутых множеств.


Предположим, что функция $\omega (Z)$ задана на множествах


${\cal D}\in \{{\cal D}\}$ и удовлетворяет условиям $1)$--$4)$.


Тогда функция ${\cal L}^p\omega (z)$ равна единице на множестве


${\cal L}^p\{{\cal E}\}$, нулю на


$\Pi^+\setminus{\cal L}^p \{{\cal D}\}$, удовлетворяет условию


Липшица в некоторой окрестности каждой точки множества


$P ={\cal L}^p\{{\cal D}\}\setminus{\cal L}^p\{{\cal E}\}$


и справедливо неравенство


\begin{equation}


\sum\limits_{\{{\cal D}\}} I(\omega,{\cal D}\setminus{\cal E})\geq 2pI({\cal


L}^p\omega,P).


\end{equation}


\end{thm}





\begin{pf}


Данное утверждение можно установить с помощью


известной техники симметризации функций [10], несколько измененной


применительно к нашему случаю. Проще, однако, воспользоваться результатами


работы [1]. Для этого запишем интегралы из неравенства (1) в виде


\begin{align*}


I(\omega,{\cal D}\setminus{\cal E})&=\lim_{\alpha\to 0}


I(\omega,{\cal D}(\alpha,\beta))+I(\omega,{\cal D}(\beta,1)),\\


%


I({\cal L}^p\omega,P)&=\lim_{\alpha\to 0}


I({\cal L}^p\omega,P(\alpha,\beta))+I({\cal L}^p\omega,P(\beta,1)),


\end{align*}


где ${\cal D}(\alpha,\beta)=\{Z\in{\cal D}:\alpha<\omega(Z)<\beta\}$,


$P(\alpha,\beta)=\{z:\alpha<{\cal L}^p\omega(z)<\beta\}=


{\cal L}^p\{{\cal D}(\alpha)\}\setminus{\cal L}^p\{\overline{{\cal


D}(\beta)}\}$, $ 0<\alpha<\beta\le 1.$


Отсюда видно, теорему 1 достаточно доказать в


предположении, что все множества ${\cal D}$ семейства $\{{\cal D}\}$ содержат


лишь конечное число точек ветвления и ограничены в совокупности конечным


числом кусочно-аналитических кривых. Проведем в полосе $\Pi$ всевозможные


прямые вида $\re z=x_k$, $k=1,\dotsc,m$, проходящие через проекции точек


ветвления множеств ${\cal D}$, ${\cal D}\in \{{\cal D}\}$,


через точки $z=0$, $z=1$, а


также прямые, являющиеся касательными к проекциям границ множеств


${\cal D}$, $0=x_1<x_2<\dotsb<x_m=1$. Присоединим к каждой связной компоненте


множества ${\cal D}$ семейства $\{{\cal D}\}$, лежащей над полосой


$\{z:x_k<x<x_{k+1}\}$, точки множества ${\cal D}$,


граничные для этой компоненты. Проекции полученных 


таким образом множеств (по всем ${\cal D}$ из $\{{\cal D}\}$)


обозначим через $D_j^k$, $j=1,\dotsc,j_k$, $k=1,\dotsc,m-1.$ Среди множеств


$D_j^k$ имеется ровно $p$ множеств, пусть это $D_{2l-1}^k$, $l=1,\dotsc,p$,


которые содержат полуполосу $\{z:x_k\leq x\leq x_{k+1},\ y\leq -c\},$


и $p$ множеств $D_{2l}^k$, $l=1,\dotsc,p$,


которые содержат полуполосу $\{z:x_k\leq x\leq x_{k+1},\ y\geq c\},$


при некотором $c>0$, едином для всех $k$. Пусть


$\wt D_{2p+t}^k=\{z:z+c_{2p+t}\in D_{2p+t}^k\}$, $t=0,1,\dotsc,j_k-2p,$ где


константы $c_{2p+t}=i|c_{2p+t}|$, $c_{2p}=0,$ подобраны так, что множества


$\wt D_{2p+t}^k$, $t=0,1,\dotsc,j_k-2p,$ попарно не пересекаются.


Введем отображения $z=\Phi_j^k(\zeta)=(x_{k+1}-x_k)(\xi+i(-1)^{j+1}\eta)+


(x_k-i(-1)^{j+1}c)$, $\zeta=\xi+i\eta\in \Pi^+_{\zeta}=


\{\zeta=\xi+i\eta:0\leq\xi\leq 1,\ \eta\geq 0\},$


и пусть $\Psi_j^k$ --- отображения,


обратные $\Phi_j^k$, $j=1,\dotsc,2p$. Рассмотрим функции


$\omega_j^k(\zeta)$, $j=1,\dotsc,2p$,


заданные в полуполосе $\Pi^+_{\zeta}$ следующим


образом:


$$


\omega_j^k(\zeta)=


\begin{cases}


1-\omega(\Phi_j^k(\zeta)), & \zeta\in \Psi_j^k(D_j^k);\\


1, & \zeta\in \Pi^+_{\zeta}\setminus\Psi_j^k(D_j^k),


\end{cases}


$$


$j=1,\dotsc,2p-1,$


$$


\omega_{2p}^k(\zeta)=


\begin{cases}


1-\omega(\Phi_{2p}^k(\zeta)+c_{2p+t}),


& \zeta\in \Psi_{2p}^k(\wt D_{2p+t}^k),t=0,


\ \dotsc,j_k-2p;\\


1, & \zeta\in \Pi^+_{\zeta}\setminus


\bigcup\limits_{t=0}^{j_k-2p}\wt D_{2p+t}^k.


\end{cases}


$$


Непосредственно из определения усредняющих преобразований видно, что в


полуполосе $\Pi^+_{\zeta}$ выполняется


$$


{\cal L}^p\omega(\Phi_{2p}^k(\zeta))=


1-{\cal L}_A(\{\omega_j^k(\zeta)\}_{j=1}^{2p}),


$$


где ${\cal L}_A$ --- линейно-усредняющее


преобразование Маркуса [1] с параметрами


$A=\{\alpha_j\}_{j=1}^{2p}$,


$\alpha_j=1/(2p)$, $j=1,\dotsc,2p.$ Из леммы 1.1 и


теоремы 1.1 работы [1] следует,


что функция ${\cal L}^p\omega(z)$ равна единице


на множестве


${\cal L}^p\{{\cal E}\}\bigcap\{z:x_k\leq x\leq x_{k+1}\},$ нулю на


$(\Pi^+\setminus {\cal L}^p\{{\cal D}\})\bigcap\{z:x_k\leq x\leq x_{k+1}\}$,


удовлетворяет условию Липшица в некоторой окрестности каждой точки множества


$P\bigcap\{z:x_k\leq x\leq x_{k+1}\}$ и справедливо неравенство


\begin{multline*}


\sum_{\{{\cal D}\}}


I(\omega,\{Z\in{\cal D}\setminus{\cal E}:


x_k<\re \pr Z<x_{k+1}\})=\sum\limits_{j=1}^{2p}


I(\omega_j^k,\{\zeta:0<\omega_j^k(\zeta)<1\})\geq \\


%


\geq 2p I({\cal L}_A\{\omega_j^k\}_{j=1}^{2p},\{\zeta: 0<{\cal


L}_A\{\omega_j^k\}_{j=1}^{2p}<1\})=


2p I({\cal L}^p\omega,P\bigcap\{z:x_k<x<x_{k+1}\}),


\end{multline*}


$k=1,\dotsc,m-1.$


Отсюда вытекает утверждение теоремы 1.


\end{pf}





\begin{note}


В условиях теоремы 1 пусть, дополнительно, множество $\Pi^+


\setminus{\cal L}^p\{{\cal E}\}$ содержится в прямоугольнике $\{z\in\Pi: 0\le


y\leq y_0\}$ при некотором $y_0>0$. Тогда из принципа Дирихле со свободной


границей имеем


\begin{equation}


I({\cal L}^p\omega,P)\geq I(y/y_0,\{z\in \Pi:0<y<y_0\})=1/y_0.


\end{equation}


Это соотношение часто бывает полезным в сочетании с теоремой 1.


\end{note}





\begin{note}


Условие $p$-листности областей совокупности $\{{\cal D}\}$ является


существенным только в контексте данного приема усреднения. Один из способов


усреднения без этого условия изучен в [11].


\end{note}





Далее рассмотрим приложения линейно-усредняющего преобразования в


исследовании свойств аналитических функций разных классов.





\section{Теоремы покрытия отрезков}





Пусть функция $w=f(z)$ регулярна в круге $U=\{z:|z|<1\}$, и пусть ${\cal R}_f$


--- риманова поверхность обратного отображения, лежащая над $w$-плоскостью


[9]. Кривую $\gamma$ на поверхности ${\cal R}_f$ назовем отрезком, если


отображение проектирования осуществляет взаимно однозначное соответствие между


$\gamma$ и некоторым отрезком $w$-плоскости. Под длиной $\gamma$ понимается


длина соответствующего отрезка. Обозначим через $\Lambda_f(\varphi)$ верхнюю


грань длин отрезков на поверхности ${\cal R}_f$, исходящих из точки $W_0=f(0)$


и лежащих над лучом $\arg w =\varphi$.





\begin{thm}


Если функция $w=f(z)=z^p+\dotsb$ регулярна и $p$-листна


в круге $U$, то для любого действительного числа $\theta$ и натурального $n$


справедливо неравенство


$$


\prod \limits_{k=1}^{n}\Lambda_f(\theta+2\pi k/n)\geq 1/4.


$$


Равенство достигается для функций


$f(z)=z^p[1+(e^{-i\theta}z^p)^n]^{-2/n}$, $p$-кратно покрывающих $w$-плоскость


с разрезами вдоль лучей


$\{w:\arg w^n = \theta n,\ |w^n| \geq 1/4 \}$.


\end{thm}





\begin{pf}


Можно считать $\theta=0$. Введем обозначения:


$\lambda_k=\Lambda_f(2\pi k/n)$,


$G_k=\{w:\mid\arg w-2\pi k/n|<\pi/n\}\setminus \{w=t\exp (2\pi ik/n):0\leq


t\leq \lambda_k\},$


$\wt G_k$ --- область $G_k$ с присоединенными к ней достижимыми граничными


точками этой области, $z=F_k(w)$ --- функция, заданная соотношениями


$$


z=(-i/\pi)\log [(\zeta-\lambda_k^{n/2})/(\zeta+\lambda_k^{n/2})],\quad


\zeta=i(w^n-\lambda_k^n)^{1/2},\quad w\in G_k,


$$


где ветви многозначных функций подобраны так, что $F_k(w)$ отображает область


$G_k$ конформно и однолистно на внутренность полосы $\Pi$, причем достижимые


граничные точки с носителем в начале переходят в $\pm i\infty$,


$k=1,\dotsc, n$.





При подходящем $\rho$, $0< \rho<1$, функция $w=f_{\rho}(z)\equiv f(\rho z)$


отображает круг $U$ на $p$-листную риманову поверхность ${\cal R}_f(\rho)$,


ограниченную аналитической кривой и имеющую лишь конечное число точек


ветвления. Сужение этой поверхности, лежащее над множеством $\wt G_k$,


образует некоторое семейство открытых множеств $\{\wt {\cal D}\}_k$ (в


относительной топологии). Функция $F_k(w)$ индуцирует соответствующее ему


допустимое семейство открытых множеств $\{{\cal D}\}_k$, лежащих над полосой


$\Pi$. Отображение множеств $\wt {\cal D}$ семейства


$\{\wt {\cal D}\}_k$ на ${\cal D}\in \{{\cal D}\}_k$ будем обозначать той


же буквой $Z=F_k(W)$. Из всех условий допустимости семейства $\{{\cal D}\}_k$


требует пояснений лишь последнее. Если некоторое множество


${\cal D}\in \{{\cal D}\}_k$ содержит дугу, целиком покрывающую прямую вида


$\re z=x$, $0<x<1$, то поверхность ${\cal R}_f(\rho)$


содержит замкнутую кривую $\gamma$, лежащую над


$\{w: \re F_k(w)=x\}$. Пусть $W_k$ --- точка кривой


$\gamma$, для которой $\arg \pr W_k=2\pi k/n$. Двигаясь от точки $W_k$ по


поверхности ${\cal R}_f(\rho)$ над лучом $\arg w=2\pi k/n$ в направлении к


$W_0$, обязательно попадем в граничную точку ${\cal R}_f(\rho)$ (по


определению $\Lambda_f(2\pi k/n)$). Двигаясь теперь из точки $W_k$ в


противоположном направлении, вновь попадем на границу, т.\,к. точка $w=\infty$


не покрывается поверхностью ${\cal R}_f(\rho)$. Таким образом, по разные


стороны от $\gamma$ имеются точки границы поверхности ${\cal R}_f(\rho)$,


что противоречит односвязности этой поверхности. Итак, семейство


$\{{\cal D}\}_k$ допустимо. Аналогично определяется соответствующее семейство


замкнутых множеств $\{{\cal E}(r)\}_k$, порожденное образом круга


$|z|\leq r$, $r<1$ (вместо $U$ в случае $\{{\cal D}\}_k$), при отображении


$w=f_{\rho}(z)$. Заметим, что множество


$\Pi^+\setminus {\cal L}^p\{{\cal E}(r)\}_k$ содержится в прямоугольнике


$\{z\in \Pi: 0\leq y\leq(-np/\pi)\log (r\rho)+(1/\pi)\log (4\lambda_k^n)+


o(1)\}$ при $r\to 0$. На множествах семейства $\{{\cal D}\}_k$ определим


функцию $\omega_k(Z)$, равную единице на ${\cal E}(r)\in \{{\cal E}(r)\}_k$,


и


$$


\omega_k(Z)=(\log|f_\rho^{-1}(F_k^{-1}(Z))|)/ \log r,\


Z\in{\cal D}\setminus


{\cal E}(r),\ {\cal D}\in \{{\cal D}\}_k.


$$


Легко видеть, что функция


$\omega_k(Z)$ удовлетворяет условиям 1)--4). Из конформной инвариантности


интеграла Дирихле следует


$$


-\frac{2\pi}{\log r}= I((\log |z|)/ \log r,\ \{z:


r<|z|<1\})=\sum\limits_{k=1}^{n}\sum_{\{{\cal D}\}_k}I(\omega_k,{\cal


D}\setminus {\cal E}(r)).


$$


С другой стороны, по теореме 1 с учетом (2) имеем


\begin{multline*}


\sum_{\{{\cal D}\}_k}I(\omega_k,{\cal D}\setminus{\cal E}(r))\geq


2p/[(-np/\pi)\log (r\rho)+(1/\pi)\log(4\lambda_k^n)+o(1)]= \\


%


=-\frac{2\pi}{\log r}\bigg[\frac{1}{n}-\frac{\log\rho}{n\log r}+


\frac{1}{np}\frac{\log(4^{1/n}\lambda_k)}{\log r}+


o\bigg(\frac{1}{\log r}\bigg)\bigg],\quad k=1,\dotsc, n.


\end{multline*}


В итоге получаем


$$


0\leq -\log\rho+


\frac{1}{np}\log\bigg(4\prod_{k=1}^{n}\lambda_k\bigg)+o(1).


$$


Переходя здесь к пределу при $r\to 0$, а затем $\rho\to 1$, завершаем


доказательство неравенства теоремы 2. Точность данной оценки проверяется


непосредственно.


\end{pf}





Полученный результат включает в себя пп.\,1--3 известной


теоремы Г.М.\,Голузина ([12], теорема 8) (без утверждения


о знаке равенства). Одновременно он содержит


многие известные теоремы для однолистных в круге функций ([4], с.\,56--57).


Заметим, что наш подход позволяет установить единственность экстремальных


отображений, а также распространить полученный результат на случай функций,


заданных в кольце. Приведем в качестве примера одну из теорем такого рода.





\begin{thm}


Пусть функция $w=f(z)$ регулярна и $p$-листна в кольце


$1\leq |z|<R$, причем


$\int\limits_{|z|=1} d\arg f(z)=2p\pi;$\ $|f(z)|\geq 1$ при


$1\leq|z|<R$ и $|f(z)|=1$,


когда $|z|=1$. Обозначим через $L_f(\varphi)$ верхнюю


грань длин отрезков на римановой поверхности обратной функции, лежащих над


лучом $\arg w=\varphi$ и содержащих на конце точку над окружностью $|w|=1$.





Пусть функция $w=E(\zeta;n,p,R)$ конформно и однолистно отображает кольцо


$1<|\zeta|<R^p$ на внешность круга $|w|>1$ с разрезами вдоль лучей $\{w:\arg


w^n=0$, $|w^n|\geq\gamma(n,p,R)\}$, $E(1;n,p,R)=1.$ Тогда для любого


действительного числа $\theta$ выполняется неравенство


$$


\prod\limits_{k=1}^{n}(L_f(\theta+2\pi k/n)+1)\geq\gamma(n,p,R).


$$


Равенство достигается для функций


$f(z)=e^{i\theta}E(\alpha z^p;n,p,R)$, $|\alpha|=1.$


\end{thm}





Доказательство этого утверждения здесь не приводится, т.\,к. оно аналогично


доказательству теоремы 2 и, кроме того, в несколько измененном варианте дано в


[11].





\section{Задача об экстремальном разбиении}





Задачи с указанным названием имеют богатую историю [13]. Рассмотрим одну из


них в обобщенной постановке. Функции $w=f_k(z)$, $k=1,\dotsc,n$, назовем


совместно $p$-листными в круге $U$, если любое значение $w$ принимается этими


функциями в совокупности не более, чем в $p$ точках с учетом кратности.


\medskip





{\bf Задача.} Определить верхнюю грань произведения


\begin{equation}


\prod\limits_{k=1}^{n}|c_k|


\end{equation}


по всем функциям $w=f_k(z)=a_k+c_kz^p+\dotsb$,


совместно $p$-листным регулярным


в круге $U$ и таким, что $|a_k|=1$, $k=1,\dotsc,n$ ($n\geq 2$).


\medskip





При $p=1$ эта задача совпадает с задачей ([13], с.\,26)


(см. также [4], с.\,53).





\begin{thm}


Максимум произведения $(3)$ равен $(4/n)^n$ и достигается


для $n$ ветвей функции $w=[(1+z^p)/(1-z^p)]^{2/n},$ каждая из которых


осуществляет $p$-кратное накрытие кругом $U$ одного из углов $\{w:|\arg w-2\pi


k/n|<\pi/n\}$, $k=1,\dotsc,n.$


\end{thm}





\begin{pf}


Введем обозначения $\theta_k=\arg a_k$, $k=1,\dotsc,n$,


$\theta_{n+1}=\theta_1+2\pi$, $\varphi_k=\theta_{k+1}-\theta_k$,


$k=1,\dotsc,n.$ Можно


считать, что $0=\theta_1<\theta_2<\dotsb<\theta_n<2\pi.$ Положим


$G_k=\{w=te^{i\theta}:0<t<\infty$, $\theta_k<\theta<\theta_{k+1}\},$ и пусть


функция $z=F_k(w)$ задана соотношениями


$z=(-i/\pi)\log[(\zeta-1)/(\zeta+1)]$,


$\zeta=(we^{-i\theta_k})^{\pi/\varphi_k}$, $w\in G_k,$


где ветви многозначных функций подобраны так, что $F_k(w)$ конформно и


однолистно отображает область $G_k$ на внутренность полосы $\Pi$, при этом


граничные точки $a_k$, $a_{k+1}$ переходят соответственно в


$+i\infty$ и $-i\infty$. При подходящем $\rho$, $0<\rho<1$, функции


$w=f_{k\rho}(z)\equiv f_k(\rho z)$


отображают круг $U$ на совместно $p$-листные


римановы поверхности ${\cal R}_{f_k}(\rho)$, лежащие над $w$-плоскостью,


ограниченные аналитическими кривыми и имеющие лишь конечное число точек


ветвления. Точку поверхности ${\cal R}_{f_k}(\rho)$, соответствующую $z=0$,


обозначим через $W_k$, $\pr W_k=a_k$, $k=1,\dotsc,n.$


Сужение этих поверхностей,


лежащее над $\overline{G}_k$, образует семейство открытых множеств


$\{\wt{\cal D}\}_k$ (в относительной топологии). Функция $F_k(w)$


индуцирует допустимое семейство открытых множеств


$\{{\cal D}\}_k$, лежащих над полосой $\Pi$ (cоответствующее отображение


множеств $\wt{\cal D}$ семейства $\{\wt{\cal D}\}_k$ на


${\cal D}\in \{{\cal D}\}_k$ обозначаем той же буквой $Z=F_k(W))$.


Заметим, что точки $W_k$ имеют окрестности, $p$-кратно покрывающие окрестности


точек $a_k$ так, что условия допустимости областей $\{{\cal D}\}_k$


легко проверяются. Аналогично определяется соответствующее семейство замкнутых


множеств $\{{\cal E}(r)\}_k$, порожденное образом круга $|z|\leq r$, $r<1$,


при отображениях $w=f_{k\rho}(z)$, $k=1,\dotsc,n.$ Простые вычисления


показывают, что множество $\Pi^+\setminus{\cal L}^p\{{\cal E}(r)\}_k$


содержится в прямоугольнике


$\{z\in\Pi: 0\leq y\leq(-1/\pi)\log[(r\rho)^p|c_k|\pi/(2\varphi_k)]+o(1)\}$,


$r\to 0.$


На множествах семейства $\{{\cal D}\}_k$ определим функцию


$\omega_k(Z)$, равную единице на ${\cal E}(r)\in\{{\cal E}(r)\}_k$, и


$$


\omega_k(Z)=(\log |f_{j\rho}^{-1}(F_k^{-1}(Z))|)/\log r,\quad


Z\in{\cal D}\setminus{\cal E}(r),\ \; {\cal D}\in\{{\cal D}\}_k,


$$


где индекс $j$ зависит от точки $Z$ и равен номеру той поверхности


${\cal R}_{f_j}(\rho)$, которой принадлежит точка $F_k^{-1}(Z)$. Как и при


доказательстве теоремы 2, конформная инвариантность интеграла Дирихле дает


$$


-\frac{2\pi n}{\log r}=\sum\limits_{k=1}^{n}\sum\limits_{\{{\cal


D}\}_k}I(\omega_k,{\cal D}\setminus{\cal E}(r)).


$$


С другой стороны, по теореме 1 с учетом (2) имеем


\begin{multline*}


\sum\limits_{\{{\cal D}\}_k}I(\omega_k,{\cal D}\setminus{\cal E}(r))\geq


2p\{(-1/\pi)\log [(r\rho)^p|c_k|\pi/(2\varphi_k)]+o(1)\}^{-1}= \\


%


=-\frac{2\pi}{\log r}\bigg[1-\frac{\log \rho}{\log r}-\frac{\log


|c_k\pi/(2\varphi_k)|}{p\log r}+o\bigg(\frac{1}{\log r}\bigg)\bigg],


\quad k=1,\dotsc,n.


\end{multline*}


Суммируя выписанные соотношения, получаем


$$


0\leq-n\log \rho-(1/p)\log\prod\limits_{k=1}^{n}|c_k\pi/(2\varphi_k)|+o(1),


\quad r\to 0.


$$


Отсюда после предельных переходов $r\to 0$, $\rho\to 1$ следует


$$


\prod\limits_{k=1}^{n}|c_k\pi/(2\varphi_k)|\leq 1.


$$


Осталось воспользоваться неравенством между средним геометрическим и средним


арифметическим


$$


\prod\limits_{k=1}^{n}\varphi_k\leq


\bigg(\,\sum\limits_{k=1}^{n}(\varphi_k/n)\bigg)^n=(2\pi/n)^n.


$$


Экстремальность указанных в формулировке теоремы функций проверяется


непосредственно.


\end{pf}





Теорема 4 обобщает известные теоремы М.А.\,Лаврентьева


$(n=2)$ и Г.М.\,Голузина


$(n=3)$ о неналегающих областях на случай $p$-листных отображений.





\section{Неравенство для рациональных функций}





Ниже предлагается новый способ получения неравенств для полиномов и


рациональных функций (ср., напр., [14]). Пусть $R(z)$ --- рациональная


функция степени $p\geq 1$, $a$ и $b$


--- комплексные числа. Предположим, что прямая


или окружность $\gamma$ отделяет корни уравнения $R(z)=a$ от корней $R(z)=b$.


Поставим задачу изучения свойств рациональной функции $R(z)$, которые вытекают


из данного ограничения. Достаточно рассмотреть случай, когда $\gamma$ ---


действительная ось, $a=0$ и $b=1$.





\begin{thm}


Пусть $w=R(z)$ --- рациональная функция степени


$p\geq 1;$ пусть $\xi_s$, $s=1,\dotsc,n$, --- корни уравнения $R(z)=0$


кратности соответственно $k_s;$ $\eta_s$, $s=1,\dotsc,m$, --- корни уравнения


$R(z)=1$


кратности $l_s$, $\sum\limits_{s=1}^{n}k_s=\sum\limits_{s=1}^{m}l_s=p$.


Предположим, что точки $\xi_s$ лежат в верхней полуплоскости, а точки $\eta_s$


--- в нижней, и пусть


\begin{alignat*}{2}


c_s&=\lim_{z\to\xi_s}|R(z)|\,|z-\xi_s|^{-k_s}, &\quad s&=1,\dotsc,n,\\


%


d_s&=\lim_{z\to\eta_s}|R(z)-1|\,|z-\eta_s|^{-l_s}, &\quad s&=1,\dotsc,m.


\end{alignat*}


Тогда имеет место точная оценка


\begin{equation}


\prod\limits_{s=1}^{n}c_s^{k_s}\prod\limits_{s=1}^{m}d_s^{l_s}\leq\frac


{\prod\limits_{s\neq t}^{n}|\xi_s-\xi_t|^{k_sk_t}\prod\limits_{s\neq


t}^{m}|\eta_s-\eta_t|^{l_sl_t}}


{\prod\limits_{s, t}^{n}|\xi_s-


\overline{\xi}_t|^{k_sk_t}\prod\limits_{s, t}^{m}


|\eta_s-\overline{\eta}_t|^{l_sl_t}}.


\end{equation}


Равенство достигается для функции


$w=(z-ih)^p[(z-ih)^p+(z+ih)^p]^{-1},$


$h>0$.


\end{thm}





\begin{pf}


Обозначим через $G$ плоскость $w$, разрезанную вдоль


отрезка $[0,1]$, и пусть $\wt G$ --- область $G$ с присоединенными к ней


достижимыми граничными точками. Пусть $z=F(w)$ --- та ветвь функции


$(-i/(2\pi))\log[w/(1-w)]$, которая конформно и однолистно отображает область


$G$ на внутренность полосы $\Pi$ так, что $F(0)=i\infty$, $F(1)=-i\infty$.


Функция $w=R(z)$ отображает верхнюю и нижнюю полуплоскости на две


непересекающиeся совместно $p$-листные римановы поверхности, лежащие над


$w$-плоскостью, ограниченные кусочно-аналитической кривой и имеющие лишь


конечное число точек ветвления. Сужение этих поверхностей, лежащее над


$\wt G$, образует семейство открытых множеств, которое переходит с помощью


отображения $F(w)$ в допустимое семейство открытых множеств $\{{\cal D}\}$,


лежащих над


полосой $\Pi$. Как и при доказательстве предыдущих теорем, отображение


римановых областей обозначаем той же буквой. Аналогично определяется


соответствующее семейство замкнутых множеств $\{{\cal E}(r)\}$, порожденное


образами кругов $|z-\xi_s|\leq r^{1/k_s}$, $s=1,\dotsc,n$, $|z-\eta_s|\leq


r^{1/l_s}$, $s=1,\dotsc,m$. Легко видеть, что множество $\Pi^+\setminus


{\cal L}^p\{{\cal E}(r)\}$ содержится в прямоугольнике $\{z\in\Pi:


0\leq y\leq(-1/(2\pi))\log r+(-1(4\pi p))\log A+o(1)\}$, $r\to0$, где


$A=\prod\limits_{s=1}^{n}c_s^{k_s}\prod\limits_{s=1}^{m}d_s^{l_s}$.


Обозначим теперь через $u(z)$ функцию, непрерывную в плоскости $z$,


равную нулю на действительной оси, единице на указанных выше кругах с центрами


в $\xi_s$ и $\eta_s$, и гармоническую в остальной части плоскости $z$ (здесь


$r>0$ достаточно мало). Из теоремы 1 работы [15] следует, что


\begin{alignat*}{2}


I(u,\{z: y>0\})&=\bigg[-\frac{1}{2\pi p}\log r+M^{+}+o(1)\bigg]^{-1}=


-\frac{2\pi p}{\log r}\bigg[1+\frac{2\pi p M^{+}}{\log r}


+o\bigg(\frac{1}{\log r}\bigg)\bigg],&\quad r&\to 0,\\


%


I(u,\{z: y<0\})&=\bigg[- \frac{1}{2\pi p}\log r+M^{-}+o(1)\bigg]^{-1}=


-\frac{2\pi p}{\log r}\bigg[1+\frac{2\pi pM^{-}}{\log r}


+o\bigg(\frac{1}{\log r}\bigg)\bigg],&\quad r&\to 0,


\end{alignat*}


где


\begin{align*}


M^{+}&=\frac{1}{2\pi p^2}\bigg\{\sum\limits_{s,t}^{n}


\log |\xi_s-\overline{\xi}_t|^{k_{s} k_{t}}-\sum\limits_{s\neq t}^{n}


\log |\xi_s-\xi_t|^{k_{s} k_{t}}\bigg\}, \\


%


M^{-}&=\frac{1}{2\pi p^2}\bigg\{\sum\limits_{s,t}^{m}\log


|\eta_s-\overline{\eta}_t|^{l_{s}l_{t}}-\sum\limits_{s\neq t}^{m}\log


|\eta_s-\eta_t|^{l_{s}l_{t}}\bigg\}.


\end{align*}


На множествах семейства $\{{\cal D}\}$ определим функцию $\omega(Z)$ по


формуле


$\omega(Z)=u(R^{-1}(F^{-1}(Z)))$, $Z\in{\cal D}$, ${\cal D}\in \{{\cal D}\}$.


Из конформной инвариантности интеграла Дирихле, теоремы 1 и неравенства (2)


имеем последовательно


\begin{multline*}


I(u,\{z: y>0\})+I(u,\{z: y<0\})=\sum\limits_{\{{\cal D}\}}


I(\omega, {\cal D}\setminus{\cal E}(r))\geq \\


%


\geq 2p[(-1/(2\pi))\log r+(-1/(4\pi p))\log A+o(1)]^{-1}= 


-\frac{4\pi p}{\log r}\bigg[1-\frac{1}{2p}\frac{\log A}{\log r}+


o\bigg(\frac{1}{\log r}\bigg)\bigg].


\end{multline*}


Суммируя полученные соотношения и переходя к пределу при $r\to 0$, приходим к


неравенству


$$


\pi p(M^{+}+M^{-})\leq -\frac{1}{2p}\log A,


$$


которое равносильно (4). Точность оценки проверяется непосредственно.


\end{pf}





Отметим частный случай неравенства (4), когда все нули функций $R(z)$ и


$R(z)-1$ первого порядка,


$$


\prod\limits_{s=1}^{p}|R'(\xi_{s})R'(\eta_{s})|\leq


\frac{\prod\limits_{s\neq t}^{p}|\xi_s-\xi_t|\prod


\limits_{s\neq t}^{p}|\eta_{s}-\eta_{t}|}


{\prod\limits_{s,t}^{p}|\xi_{s}-\overline{\xi}_{t}|\prod\limits_{s,t}^{p}


|\eta_{s}-\overline{\eta}_{t}|}


<2^{-2p}\prod\limits_{s=1}^{p}|\im \xi_{s} \im \eta_{s}|^{-1}.


$$


Знак равенства в первом неравенстве достигается, например, для функции


$R(z)=\frac{9}{80}\big[\big(\frac{z-i}{z+i}\big)^2-\frac{1}{9}\big]$


($p=2$).


Описание всех случаев равенства здесь не приводится, т.\,к.


требует привлечения техники обобщенных приведенных


модулей [15] и теорем единственности для усредняющих преобразований [2], что


существенно увеличит объем статьи. Полагая


$R(z)=P(z)\equiv\prod\limits_{s=1}^{p}(z-\xi_s)$,


получим оценку для многочленов


$\prod\limits_{s, t}^{p}|\xi_s-\overline{\xi}_t|


\prod\limits_{s, t}^{p}|\eta_s-\overline{\eta}_t|\leq 1.$


Данная оценка является, по-видимому, неточной в классе всех многочленов


$P(z)=z^p+c_{p-1}z^{p-1}+\dotsb+c_0$, удовлетворяющих указанным выше условиям.





\begin{thebibliography}{99}





\bibitem{1}


Marcus M. {\em Radial averaging of domains, estimates for Dirichlet integrals


and applications} // J. Anal. Math. -- 1974. -- V.27. -- P.47--78.


\bibitem{2}


Митюк И.П. {\em Симметризационные методы и их применение в геометрической


теории функций. Введение в симметризационные методы}. -- Краснодар: Изд-во


Кубанск. гос. ун-та, 1980. -- 91~с.


\bibitem{3}


Митюк И.П. {\em Применение симметризационных методов в геометрической теории


функций}. -- Краснодар: Изд-во Кубанск. гос. ун-та, 1985. -- 94 с.


\bibitem{4}


Дубинин В.Н. {\em Симметризация в теории функций комплексного переменного} //


УМН. -- 1994. -- Т.\,49. -- Вып.\,1. -- С.\,3--76.


\bibitem{5}


Дженкинс Дж. {\em Однолистные функции и конформные отображения}.


-- М.: Ин. лит., 1962. -- 265 с.


\bibitem{6}


Garabedian P.R., Royden H.L. {\em The one-quarter theorem for mean univalent


functins} // Ann. Math. -- 1954. -- V.\,59. -- \No\,2. -- P.\,316--324.


\bibitem{7}


Krzyz J. {\em Distortion theorems for bounded $p$-valent functions} //


Ann. Univ. M. Curie-Sklodowska. Sec.\,A. -- 1958. -- V.\,12. -- P.\,29--38.


\bibitem{8}


Шлык В.А. {\em Некоторые оценки в кольце для слабо однолистных функций,


выпускающих значения на окружности} // Изв. вузов. Математика. -- 1981.


-- \No\,8. -- С.\,85--86.


\bibitem{9}


Стойлов С. {\em Теория функций комплексного переменного}. Т.\,2.


-- М.: Ин. лит., 1962. -- 416 c.


\bibitem{10}


Хейман В.К. {\em Многолистные функции}. -- М.: Ин. лит., 1960. -- 180 c.


\bibitem{11}


Дубинин В.Н. {\em Некоторые применения линейно-усредняющего преобразования в


теоремах покрытия}. -- Кубанск. гос. ун-т. -- Краснодар, 1977. -- 22 c. --


Деп. в ВИНИТИ 28.02.77, \No\,770--77.


\bibitem{12}


Голузин Г.М. {\em Некоторые теоремы покрытия в теории аналитических функций}


// Матем. сб. -- 1948. -- Т.\,22. -- \No\,3. -- С.\,353--372.


\bibitem{13}


Кузьмина Г.В. {\em Методы геометрической теории функций}. II // Алгебра и


анализ. -- 1997. -- Т.\,9. -- Bып.\,5. -- С.\,1--50.


\bibitem{14}


Borwein P., Erdelyi T. {\em Polynomials and polynomial inequalities}. --


New York: Springer-Verlag, 1995. -- 480 p.


\bibitem{15}


Дубинин В.Н. {\em Асимптотика модуля вырождающегося конденсатора и


некоторые ее применения} // Зап. научн. семин. ПОМИ. -- 1997. -- Т.\,237.


-- С.\,56--73.





\end{thebibliography}





\vskip 2cm





\post{\it Институт прикладной математики }{\it Поступила}


\post{\it Дальневосточного отделения}{18.01.1999}


\post{\it Российской Академии наук}{}





\label{end}


\end{document}


